\chapter{Rebuilding Value in the Shear Zone}

When a capability becomes cheap to obtain, the social position once attached to it can move. The destination depends on demand, ownership, bargaining power, and public institutions. Three things must be kept distinct: a function's social necessity, the reward given to whoever performs it, and the dignity of the person involved.

Verification, care, sustained relationships, and responsibility can become relatively scarce in a world of abundant generation. Turning that scarcity into support and recognition requires arrangements that can pay for it.

\section{A History of Capabilities Leaving the Body}

\subsection{Writing Changed What Memory Proved}

In Plato's account, Socrates worries that writing can supply the appearance of wisdom without its internal possession. A stored statement can be consulted without the learner acquiring the judgment behind it.

The concern identifies a recurring distinction. External storage changes what a person's immediate performance demonstrates. It also creates new kinds of thought: a written proof can be paused, checked, and challenged by someone who never met its author.

Oral tradition itself was more inventive than a storage metaphor suggests. Parry and Lord's work on epic singers described composition through reusable patterns and formulas, with performances varying rather than reproducing one fixed text \cite{parry1971making,parry1960singer}. Writing changes a practice of memory and composition, not merely the location of a recording.

Modern experiments on expected access to information show another reallocation: remembering where to retrieve material can improve while memory for the material changes \cite{sparrow2011google}. External memory can support a stronger system when the person retains enough internal knowledge to recognize relevance, error, and missing context.

The migration is from holding every item to knowing how to use and assess what remains accessible. Access itself still depends on records, preservation, tools, and institutions.

\subsection{Cheap Copies Changed Authority}

Printing lowered the cost of reproducing texts and helped reorganize correction, comparison, distribution, and controversy \cite{eisenstein1979printing}. It expanded audiences and created work around editions, interpretation, and trust.

The history is uneven across regions and technologies. Movable type in Europe was one important transformation within a much longer history of writing and printing. The point is the changing relation between possession and authority.

When access to a text becomes easier, ownership of a copy supplies less distinction. A dependable edition, a persuasive interpretation, or a trusted publisher can become more valuable. Cheap copying can spread errors along with corrections.

Computer systems extended the movement to rule execution. Calculation, sorting, accounts, retrieval, and repeatable workflows moved into machinery. A professional's position could shift toward choosing the model, interpreting exceptions, checking the result, and assuming a defined responsibility.

Such shifts create winners, losses, and new dependencies. They do not conserve the number of jobs or the value of every old skill. Institutions help determine whether the increased productive capacity becomes broadly shared benefit.

\subsection{Clocks and Engines Reorganized the Operator}

A mechanical clock made intervals publicly comparable and helped support schedules that could be imposed across workers and machines. The resulting coordination expanded productive possibilities while changing control over the working day.

Steam power moved large quantities of work into fuel, engines, factories, and transport networks \cite{landes1969unbound}. Human muscle remained necessary in many tasks, yet it was no longer the only scalable source of industrial power.

The operator's problem changed. Keeping machines, fuel, schedules, and people coordinated could become more important than supplying the direct force. The new positions depended on access to capital, education, authority, and organization.

These histories expose a repeated pattern: technology changes which evidence a society uses to recognize competence, and which institution controls the next scarce interface.

\subsection{Formal Discovery Is Moving Too}

Machine learning has reached language, code, design, scientific hypotheses, and mathematical search. The artifact once treated as evidence of years of preparation can now be produced through a different route.

Mathematics is a demanding case because discovery and checking can be distinguished sharply. AlphaGeometry's 2024 work combined neural proposals with symbolic deduction \cite{trinh2024alphageometry}. In May 2026, an OpenAI-generated construction challenged the conjectured near-linear exponent for the planar unit-distance problem; a group of mathematicians supplied a concise exposition and analysis \cite{alon2026remarks}.

The intellectual contribution is substantial. It changes which problems are worth asking and which techniques can be combined. The knowledge still includes the precise statement, its proof, its relation to earlier work, and the consequences other researchers can extract.

September's OpenAI announcement on Navier--Stokes went further, presenting a smooth-forcing finite-time singularity construction and a Lean formalization \cite{openai2026navierstokes}. The theorem's domain, force assumptions, and correspondence to the stated problem are central parts of evaluating it. “The model solved fluids” would discard the very content that makes the result mathematical.

The current lesson is that verification itself can accelerate. Formal tools can check derivations cheaply once an adequate statement and proof exist. The remaining human and institutional work includes specification, understanding, review, priority, pedagogy, and choosing what deserves attention.

Tao's evolving account of AI-assisted mathematics emphasizes these distinctions between candidate production, independent checking, formalization, and understanding \cite{tao2026aiviews}. The scarce interface can move again when a previously laborious checking step becomes automated.

An illustrative capacity ratio is
\begin{equation}
 \rho=\min\left\{1,\frac{d}{g}\right\},
\end{equation}
where $g$ is candidate production and $d$ comparable digestion or review capacity, in matching units. In the unsaturated region $0<d/g<1$, $\rho$ falls when $g$ grows proportionally faster than $d$. At the capped value one, the same strict derivative statement need not hold.

The quantity is useful only after defining what counts as a candidate and what digestion requires. Reviewing a formal proof, understanding it, and integrating it into a research field are different activities.

\subsection{Capability, Qualification, and Permission}

Embodied action and biological intervention extend the same change into heavier media. A generated policy needs an actuator and an operating domain. A proposed sequence needs a biological process capable of using it. A demonstrated capability needs an account of where it can be relied upon.

Traditional qualifications bundled several signals: preparation, experience, exclusivity, and responsibility. Spence's signaling framework helps explain how some qualifications communicate information under particular incentive conditions \cite{spence1973job}. Machine assistance can loosen the bundle.

The response should inspect each function. Does the assessment establish the skill needed? Does the trainee still obtain the experience needed for later judgment? Does the professional have authority matching their duty? Does the user know which part of a result was actually checked?

The history offers no automatic promise that value migrates fairly. It offers a practical method: identify the capability being externalized, the old signal losing force, and the new interface on which reliable performance depends.

The next four principles ask how those interfaces should be maintained.
